\section{Evidence Boundaries}

The chapter does not claim that neuroscience has solved consciousness, nor
that the cited papers directly measure free will. It also does not claim that
every local behavioural pattern should be mapped one-to-one onto a whole-brain
state. The \textbf{strong} conclusion is narrower: organized state structure,
stability, and transition dynamics matter biologically.

The nested-state reading is therefore \textbf{medium} evidence, because it
transfers a biologically grounded vocabulary into the book's behavioural
framework. Anything stronger, especially claims about the physical basis of
subjective awareness, would become \textbf{speculative} and should remain
outside the chapter core.

\begin{discussion}
The boundary can be stated operationally:
\begin{itemize}[nosep]
  \item \textbf{Strong}: state organization, stability structure, and
        transition dynamics are biologically serious.
  \item \textbf{Medium}: behavioural frames can be discussed using that
        vocabulary when the transfer is explicit and disciplined.
  \item \textbf{Speculative}: claims about subjective awareness itself, exact
        neural realization of ordinary tasks, or consciousness-physics
        unification remain outside the empirical chapter core.
\end{itemize}
\end{discussion}
